\documentclass[../main.tex]{subfiles}
\graphicspath{{\subfix{../images/}}}
\begin{document}
% \textbf{Hello world!}
% \begin{figure}[bh]
% \centering
% \includegraphics[width=3cm]{automata/tautomata.pdf}

% \label{fig:img1}
% \caption{Overleaf logo}
% \end{figure}

\chapter{Applications of VQE I: Chemistry}

\section{Introduction: The Why of quantum computing for quantum chemistry}

\begin{itemize}
    \item In the last chapter we discussed working principles f VQE and how such algorithm could be implemented on a measuremnt-based quantum cmoputer respecting determinism. We also gave some examples as benchmark (list what was studied in mbqc paper)
    \item Next, we want to dive deeper what applications could be given such 
    \item  check intro of paper 1
\end{itemize}

\section{Introduction to Quantum computational chemistry}



\begin{itemize}
    \item the electronic structure hamiltonian and the Fock space
    \item second quantization
    \item cholesky based decomposition of the ahmiltonian
    
    \item mapping to qubits: Jordan Wiegner
    \item connection to matchgate, fermionic Gaussian
    
    \item AFQMC fundamentals, fermionic sign problem
    \item Projector methods in imaginary time, Hubbard–Stratonovich transformation, random walkers (Slater determinants), local energy estimator, and the phase/sign problem; phaseless approximation and role of the trial state

    \item VQE for quantum chemistry
    \item Overview of variational principle, ansätze (e.g., UCC), and measurement optimization strategies (grouping, basis rotations, classical shadows) to lower measurement counts on NISQ.

    \item Classical shadows for quantum chemistry
    \item Core idea of classical shadows (randomized measurements) and why they are attractive when many observables of the same state are needed; introduce Clifford, matchgate, and orbital-rotation shadows; emphasize variance/shot-noise behavior and reuse of measurements.

    \item QC‑AFQMC (using quantum trial states in AFQMC)
    \item Where the quantum state enters: overlaps needed for importance sampling (overlap ratio), force-bias evaluation, and local energy; scaling and factorization to $O(N^3)$ for local energy with modern compression; total overlap burdens across walkers/blocks/steps; mean-field subtraction requires trial-state expectations.
    \item Practical measurement challenge and post-processing costs; with classical-shadow schemes the classical post-processing for QC‑AFQMC overlap estimation can scale as $O(N^9)$ per walker end-to-end (algorithmic workflow perspective).
    
\end{itemize}

This section aims at a complete yet minimal introduction to the electronic structure theory maianly studied in quantum chemistry in so-called second-quantized formulation with the aim to cover all concepts to makes the motivation, methods and results of the publications \todo{cite bith paper} comprehensive. This will include a minimalistic introduction to quantum chemistry and is formulation, the expectation the field has on quantum computer, and how such problems are mapped onto the quantum computer. A quick introduction on what a quantum computer is, measuremnts are done or algorithms as the VQE is given \todo{ref the intro}. For more details on electronic structure theory, we recommend the study of Refs.~\cite{helgaker2014molecular, szabo2012modern}, or more taiored towards quantum computing practitoiners Refs. \cite{mcardle2020quantum, cao2019quantum}.


\subsection{The electronic structure theory in second quantuization}





\section{Quantum‑enhanced AFQMC in industrial chemistry and materials }

\begin{itemize}
    \item write what we did: proof of concept on challenging test systems
    \item second paper, demystyfication of classical resources required run qqc-qmc. 
\end{itemize}

\subsection{Proof of concept: da toysystems}

Objective and scope (industrial view)
Test QC‑enhanced AFQMC on representative systems: H4 potential curve; relative energies of 1O2 and O3 vs 3O2; and a quasi‑1D model derived from CuBr2 mapped to a Hubbard-like Hamiltonian.



Trial states and set‑up
Classical: HF, CASSCF/CAS; quantum: UCCSD‑VQE initialised by CCSD amplitudes; also MPS trials for the lattice model; implementation via ipie, PySCF, and cirq; CBS extrapolation for molecular benchmarks and ZPVE corrections for experiment comparison.




Molecular benchmarks (H4, O2/O3)
H4: AFQMC(VQE) systematically improves VQE energies and tracks FCI closely, but does not attain chemical accuracy across all geometries.
O2/O3: AFQMC with CAS trials yields relative energies within or near chemical accuracy; caveats include AFQMC statistical error propagation, CBS extrapolation from limited basis points, and ZPVE consistency

Materials benchmark (CuBr2→1D Hubbard-like model)
Trial-state quality and symmetry matter: higher-fidelity MPS trials need not yield better AFQMC energies than symmetry-preserving HF; the structure of the trial state (symmetries) can dominate performance.

Key takeaways
For molecules, QC‑AFQMC with high-quality small active-space trials is promising against industrially relevant targets; for lattice models, enforce symmetries and physical structure in trials; integration into workflows and robust ZPVE/forces remain as practical bottlenecks to address.

\subsection{Scaling study:}

QC‑AFQMC requires many overlaps between the quantum trial state and evolving (non‑orthogonal) Slater-determinant walkers; evaluate measurement and classical post-processing costs, error propagation, and scaling.

Where overlaps appear and how many
Importance sampling and force bias require $O(N^2)$ per walker per step; local energy is periodic but dominant, $O(N^3)$ with factorization; aggregate across walkers/blocks yields substantial measurement burdens

Measurement schemes
Hadamard test: direct overlap estimation, but controlled-state preparation increases circuit depth; sampling independent of N, but impractical for NISQ depth and per‑overlap repetition.
Classical shadows:
Clifford shadows: few shots for many observables, but exponentially costly classical post‑processing for Slater‑determinant overlaps.
Matchgate shadows: polynomial classical post‑processing via Pfaffian interpolation; polynomial shot complexity; efficient for Slater determinant overlaps.
%Orbital-rotation shadows: particle‑number‑preserving ensemble; shot variance scales with particle number η and can decrease with N; needs ancillas and deeper circuits; efficient in $\eta$ with adapted interpolation methods.

Statistical behavior and covariances
Overlap estimates are near‑Gaussian and exhibit shot‑noise scaling with the number of snapshots; covariances arise when reusing shadows for correlated walker states over imaginary time and should be mitigated by resampling/reblocking techniques.


Error propagation and end‑to‑end cost
Derive how overlap errors propagate into local energy and final AFQMC energy; errors scale with Hamiltonian 2‑norm, not 1‑norm as in VQE; trading quantum shots vs. classical walkers/blocks can compensate within limits.
Overall, combining measurement and classical post‑processing leads to an $O(N^9)$ per‑walker scaling for QC‑AFQMC overlap estimation pipelines under state‑of‑the‑art schemes.
Numerical demonstrations
H4 and H2O with QNP‑VQE trials: AFQMC remains accurate up to noise thresholds (~1e‑2–1e‑3 overlap error), but full‑pipeline noise (propagation + local energy) is more damaging than propagation‑only; larger bases (cc‑pVDZ/VTZ) highlight sensitivity to overlap noise in propagation.

key takeaways 

\end{document}